\documentclass[11pt]{article}
\usepackage[a4paper,margin=2.6cm]{geometry}
\usepackage{amsmath,amssymb,amsthm,mathtools}
\usepackage{booktabs}
\usepackage[hidelinks]{hyperref}

\newtheorem{theorem}{Theorem}
\newtheorem{proposition}[theorem]{Proposition}
\newtheorem{lemma}[theorem]{Lemma}
\newtheorem{corollary}[theorem]{Corollary}
\theoremstyle{remark}
\newtheorem{remark}[theorem]{Remark}

\newcommand{\SL}{\mathrm{SL}}
\newcommand{\Mp}{\mathrm{Mp}}
\newcommand{\Z}{\mathbb{Z}}
\newcommand{\C}{\mathbb{C}}
\newcommand{\R}{\mathbb{R}}
\newcommand{\Hb}{\mathbb{H}}
\newcommand{\tr}{\operatorname{tr}}
\newcommand{\adj}{\operatorname{adj}}

\title{The theta triplet of $\Gamma(4)$: a nilpotent field, a determinant identity,\\ and a period phase along closed geodesics}
\author{origin-axiom}
\date{Draft 1, \today}

\begin{document}
\maketitle

\begin{abstract}
The squares of the three even theta constants span the weight-one modular forms of $\Gamma(4)$, and Jacobi's identity
makes the vector $\rho=(-\theta_2^2,-\theta_4^2,i\theta_3^2)$ isotropic. We record four exact statements built on it.
(i) With the three quaternion units placed at the half-periods where $\theta_2$, $\theta_4$ and $\theta_3$ vanish, a field
$\Phi_\rho$ on the torus $\C/(\Z+\tau\Z)$ is nilpotent at every point; the proof takes two lines, and the other placements
fail at a generic $\tau$. (ii) For a bilinear map $Y:M_2(\C)\times M_2(\C)\to M_3(\C)$ built from the quaternion units,
$\det Y(S,M)=4\det S\,\det M\,\tr(\epsilon SM^T)$, and the adjugate splits into two terms, each a determinant times a quadratic
form. (iii) The real-analytic function $D=-\theta_3^4h(\lambda)/(4\pi\eta^6)$, with $\lambda=\theta_2^4/\theta_3^4$ and
$h(\lambda)=\lambda\log|\lambda|+(1-\lambda)\log|1-\lambda|$, transforms with weight $-1$ and the multiplier of $\eta^6$.
(iv) Hence, along the closed geodesic of a hyperbolic element written as a positive word in the two unit shears $L$ and
$R$, the amplitude $D\,(dz/ds)^{-1/2}$ returns after one period multiplied by exactly $i^{\#L-\#R}$. We also compute the
average over the modular surface of $f=|D|/\|\rho\|^2$, $0.31500491$. Scripts for every computation are public.
\end{abstract}

\section{Introduction}

The vector $u=(\theta_2^2,\theta_3^2,\theta_4^2)$ of squared theta constants is a basis of the weight-one forms of
$\Gamma(4)$; the graded ring of $\Gamma(4)$ is $\C[u]$ modulo Jacobi's identity $\theta_3^4=\theta_2^4+\theta_4^4$. The
modular group permutes the three components through $\SL(2,\mathbb F_2)\cong S_3$, which is the action on the three even
theta characteristics of the torus. Modular flavour models use $u$ as their weight-one triplet~\cite{NovichkovPenedoPetcov2021,LiuYaoDing2020}.
This note records exact statements about $u$ that arose in a larger computational project~\cite{OriginAxiomRepo}, stated
here on their own. Computations in that public record are cited by their label, such as B1666.

\paragraph{Notation.} $\tau\in\Hb$, $q=e^{i\pi\tau}$, and $\theta_2,\theta_3,\theta_4$ are the theta constants in the nome $q$;
$\eta$ is Dedekind's eta function and $\chi_\eta$ its multiplier, a character of the metaplectic group $\Mp_2(\Z)$ of order
$24$. The unit shears are $L=\begin{psmallmatrix}1&1\\0&1\end{psmallmatrix}$ ($\tau\mapsto\tau+1$) and
$R=\begin{psmallmatrix}1&0\\1&1\end{psmallmatrix}$ ($\tau\mapsto\tau/(\tau+1)$). The torus is $E_\tau=\C/(\Z+\tau\Z)$, with
Weierstrass function $\wp$ and $e_1=\wp(\tfrac12)$, $e_2=\wp(\tfrac{1+\tau}2)$, $e_3=\wp(\tfrac\tau2)$, so that
$e_1-e_2=\pi^2\theta_4^4$, $e_2-e_3=\pi^2\theta_2^4$, $e_1-e_3=\pi^2\theta_3^4$~\cite{WhittakerWatson1927}. Set
\[
  \rho=(\rho_i,\rho_j,\rho_k)=(-\theta_2^2,\,-\theta_4^2,\,i\theta_3^2),\qquad \rho\cdot\rho=\textstyle\sum_r\rho_r^2=\theta_2^4+\theta_4^4-\theta_3^4=0 .
\]
The scale of $\rho$ plays no role below except in Section~\ref{sec:D}, where $\|\rho\|^2=\sum_r|\theta_r|^4$ up to a common
factor that cancels in $f$.

\paragraph{The multipliers.} Up to scalars the action of $\SL(2,\Z)$ on $u$ is the permutation of the even
characteristics just described, with $\theta_2^2(\tau+1)=i\theta_2^2(\tau)$ and $\theta_3^2\leftrightarrow\theta_4^2$. The genuine
lifts of this projective action to $\Mp_2(\Z)$ that agree with the multiplier $\rho_u$ of $u$ projectively are
$\chi_\eta^k\otimes\rho_u$, $k\in\Z/24$, because two genuine representations that agree projectively differ by a character
and the abelianisation of $\Mp_2(\Z)$ is $\Z/24$, generated by $\chi_\eta$. Since $\rho_u(s)^2=-1$ and $\chi_\eta(s)^2=e^{-i\pi/2}$,
a weight $w$ is compatible with the central element for $\chi_\eta^k\otimes\rho_u$ exactly when
\[
  w\equiv1+\tfrac k2\pmod 2 .
\]
A particular lift computed in~\cite{OriginAxiomRepo} (B1634) is the case $k=21$, for which weight $\tfrac32$ is
compatible; which $k$ occurs depends on a choice of lifts, and we do not use it below.

\section{A nilpotent field on the torus}\label{sec:nil}

Let $\mathsf e_i,\mathsf e_j,\mathsf e_k\in M_2(\C)$ be quaternion units: $\mathsf e_r^2=-1$ and $\mathsf e_r\mathsf e_s=-\mathsf e_s\mathsf e_r$ for
$r\neq s$. Place them at the half-periods where their theta functions vanish: $z_i=\tfrac12$ (where $\theta_2(\pi z)$
vanishes), $z_j=\tfrac\tau2$ ($\theta_4$) and $z_k=\tfrac{1+\tau}2$ ($\theta_3$). For each $r$ let $\hat f_r$ be the meromorphic
function on a double cover with $\hat f_r^{\,2}=\wp(z)-\wp(z_r)$ and residue one at $z=0$.

\begin{proposition}\label{prop:nil}
The field $\Phi_\rho=\sum_r\rho_r\hat f_r\,\mathsf e_r\,dz$ satisfies $\Phi_\rho^2=0$ at every point of $E_\tau$, for every $\tau$.
Of the six placements of the units at the three half-periods, this is, up to signs, the only one with this property at
a generic $\tau$.
\end{proposition}
\begin{proof}
The units anticommute, so the cross terms cancel and
$\Phi_\rho^2=-\sum_r\rho_r^2\hat f_r^{\,2}\,dz^2=-\bigl[(\rho\cdot\rho)\,\wp(z)-\sum_r\rho_r^2\wp(z_r)\bigr]dz^2$. The first term vanishes
by Jacobi's identity. For the second,
\[
  \textstyle\sum_r\rho_r^2\wp(z_r)=\theta_2^4e_1+\theta_4^4e_3-\theta_3^4e_2
  =\theta_2^4(e_1-e_2)-\theta_4^4(e_2-e_3)=\pi^2\bigl(\theta_2^4\theta_4^4-\theta_4^4\theta_2^4\bigr)=0,
\]
using $\theta_3^4=\theta_2^4+\theta_4^4$. For the other placements the same sum is a non-zero combination of the $e_r$ at a
generic $\tau$; the script \texttt{nilpotent.py} finds it at least $0.43$ in absolute value at three generic values of
$\tau$, and vanishing to $10^{-30}$ for the stated placement.
\end{proof}

\section{A determinant identity}\label{sec:det}

Let $\epsilon=\begin{psmallmatrix}0&1\\-1&0\end{psmallmatrix}$ and $(C_1,C_2,C_3)=(\epsilon,\ \mathrm{diag}(i,-i),\ \mathrm{diag}(i,-i)\,\epsilon)$,
and define the bilinear map $Y:M_2(\C)\times M_2(\C)\to M_3(\C)$ by
\[
  Y(S,M)_{ab}=\sum_{m,n}\bigl(C_a^{T}\epsilon\,S\,C_b\bigr)_{mn}M_{mn}.
\]
For $M=\begin{psmallmatrix}m_0&m_1\\m_2&m_3\end{psmallmatrix}$ let
\[
  A(M)=\begin{pmatrix}
  m_0^2+m_1^2+m_2^2+m_3^2 & 2i(m_0m_2+m_1m_3) & -i(m_0^2+m_1^2-m_2^2-m_3^2)\\
  2i(m_0m_1+m_2m_3) & -4m_1m_2 & 2(m_0m_1-m_2m_3)\\
  -i(m_0^2-m_1^2+m_2^2-m_3^2) & 2(m_0m_2-m_1m_3) & -m_0^2+m_1^2+m_2^2-m_3^2
  \end{pmatrix},\qquad P=\mathrm{diag}(1,-1,-1).
\]

\begin{theorem}\label{thm:det}
For all $S,M\in M_2(\C)$,
\[
  \det Y(S,M)=4\det S\,\det M\,\tr(\epsilon SM^T),\qquad \adj Y(S,M)=\det S\,A(M)+\det M\,P\,A(S).
\]
With $N=M\epsilon$, the map is symmetric, $Y(S,N\epsilon^{-1})=Y(N,S\epsilon^{-1})$, and
$\tr(\epsilon SM^T)=-\bigl(\det(S+N)-\det S-\det N\bigr)$ is minus the polar form of the determinant.
\end{theorem}
\begin{proof}
All four statements are polynomial identities in the eight entries of $S$ and $M$. The script \texttt{joint\_law.py} expands
both sides symbolically and finds every residual zero. As a control, replacing $P$ by the identity, or the factor $4$ by
$3$, breaks the corresponding identity. We have not found the identity in the literature; it may be a known identity of
the quaternion algebra in disguise.
\end{proof}

\begin{corollary}\label{cor:det}
(1) $Y(S,M)$ is singular exactly when $\det S=0$, $\det M=0$ or $\tr(\epsilon SM^T)=0$.
(2) The squares $x$ of the singular values of $Y$ are the roots of $x^3-e_1x^2+e_2x-e_3$ with $e_1=\|Y\|^2$,
$e_2=\|\det S\,A(M)+\det M\,PA(S)\|^2$ and $e_3=16|\det S\det M\tr(\epsilon SM^T)|^2$.
\end{corollary}
\begin{proof}
Both follow from the theorem, since $e_3=|\det Y|^2$ and $e_2=\|\adj Y\|^2$ (Frobenius norm) is the sum of the products of
pairs of squared singular values.
\end{proof}

\begin{remark}[an observation without proof]
On the diagonal $N=S$ the singular values are in geometric progression, $\sigma_1\sigma_3=\sigma_2^2$; equivalently, by
Theorem~\ref{thm:det}, $\sigma_2^3=|\det Y|$. The script \texttt{joint\_law.py} checks this on twenty random $S$ (numpy
generator seed $7$) to $10^{-14}$. We do not have a proof.
\end{remark}

\section{A form of weight $-1$}\label{sec:D}

Let $\lambda=\theta_2^4/\theta_3^4$, the modular lambda function, $h(\lambda)=\lambda\log|\lambda|+(1-\lambda)\log|1-\lambda|$, and
\[
  D(\tau)=-\frac{\theta_3^4\,h(\lambda)}{4\pi\,\eta^6}.
\]
It arises in~\cite{OriginAxiomRepo} (B1654, B1657) from the norms of Szeg\H{o} kernels on the punctured torus, but the
definition here stands alone.

\begin{lemma}\label{lem:auto}
$D$ has weight $-1$ and the multiplier $\chi$ of $\eta^6$: $D(g\tau)=\chi(g)(c\tau+d)^{-1}D(\tau)$ for
$g=\begin{psmallmatrix}a&b\\c&d\end{psmallmatrix}\in\SL(2,\Z)$, with $\chi(L)=i$ and $\chi(R)=-i$.
\end{lemma}
\begin{proof}
Under $\tau\mapsto\tau+1$, $\theta_3\leftrightarrow\theta_4$ and $\theta_2\mapsto e^{i\pi/4}\theta_2$, so $\lambda\mapsto\lambda/(\lambda-1)$ and
$\theta_3^4\mapsto\theta_4^4=\theta_3^4(1-\lambda)$. A short computation gives $h(\lambda/(\lambda-1))=-h(\lambda)/(1-\lambda)$, and
$\eta^6\mapsto i\eta^6$. Hence $D(\tau+1)=iD(\tau)$. Under $\tau\mapsto-1/\tau$, $\theta_3\mapsto\sqrt{-i\tau}\,\theta_3$ and
$\theta_2\leftrightarrow\theta_4$ with the same factor, so $\lambda\mapsto1-\lambda$, $\theta_3^4\mapsto-\tau^2\theta_3^4$, $h(1-\lambda)=h(\lambda)$,
and $\eta^6\mapsto(-i\tau)^3\eta^6=i\tau^3\eta^6$. Hence $D(-1/\tau)=i\tau^{-1}D(\tau)$. Both sides of the law compose by the
cocycle rule and the two maps generate the group, so the law holds throughout, with $\chi$ well defined because $D$ is
not identically zero; as a consistency check, $\chi(S_0)=\chi(L)=i$ and $(S_0L)^3=-I$ give $\chi(-I)=i^6=-1=i^2=\chi(S_0^2)$. Then $\chi(L)=i$, and
$R=S_0L^{-1}S_0^{-1}$ with $S_0=\begin{psmallmatrix}0&-1\\1&0\end{psmallmatrix}$ gives $\chi(R)=\chi(L)^{-1}=-i$.
\end{proof}

$D$ is real-analytic and not holomorphic, because of $\log|\cdot|$; the law is nonetheless exact. With
$\|\rho\|^2=\sum_r|\theta_r|^4$, the function
\[
  f=\frac{|D|}{\|\rho\|^2/(4\pi^2|\eta|^6)}=\frac{\pi\,|\theta_3|^4\,|h(\lambda)|}{|\theta_2|^4+|\theta_3|^4+|\theta_4|^4}
\]
is invariant under $\SL(2,\Z)$. It vanishes at $\omega=e^{2\pi i/3}$, where $\lambda=e^{\pm i\pi/3}$ and $h(\lambda)=0$, and toward the
cusp, where $f\sim|q|\,\bigl|8\pi^2y+4\pi(1+e^{-2\pi ix})-32\pi\ln2\bigr|$ for $\tau=x+iy$, $y\to\infty$ (the script
\texttt{cusp\_law.py} of~\cite{OriginAxiomRepo} checks the leading term to $1.2\cdot10^{-6}$ at $y=5$). Its value at $i$ is
$1.0888$.

\section{The period phase}\label{sec:phase}

Let $\gamma$ be a hyperbolic element of $\SL(2,\Z)$ written as a positive word in $L$ and $R$ (every hyperbolic conjugacy
class has such a representative up to sign~\cite{Series1985}), and let $z(s)$, $s\in\R$, be its axis in $\Hb$ parametrised by
hyperbolic arclength from the repelling to the attracting fixed point, so that $z(s+\ell)=\gamma z(s)$ with $\ell$ the
translation length. Define the amplitude $\psi(s)=D(z(s))\,(dz/ds)^{-1/2}$, the square root continued along the axis.

\begin{theorem}\label{thm:phase}
$\psi(s+\ell)=i^{\#L-\#R}\,\psi(s)$ for every $s$, where $\#L$ and $\#R$ count the letters of the word.
\end{theorem}
\begin{proof}
Write $\gamma=\begin{psmallmatrix}a&b\\c&d\end{psmallmatrix}$. Since $z(s+\ell)=\gamma z(s)$, $dz/ds$ at $s+\ell$ is $(cz+d)^{-2}$ times its
value at $s$, and $D(z(s+\ell))=\chi(\gamma)(cz+d)^{-1}D(z(s))$ by Lemma~\ref{lem:auto}. Let $r(s)$ be the ratio of the continued
square root $(dz/ds)^{-1/2}$ at $s+\ell$ to its value at $s$. Then $r^2=(cz+d)^2$, and since $cz+d$ has no zero on $\Hb$ and $r$
is continuous, $r=\epsilon\,(cz+d)$ with a single sign $\epsilon$. As $s\to-\infty$ both $z(s)$ and $z(s+\ell)$ tend to the
repelling fixed point $x$ and the tangent directions at the two points agree; for $s$ negative enough the argument of
$dz/ds$ varies by less than $\pi$ along $[s,s+\ell]$, so the continued square root takes the principal branch of a ratio
that tends to a positive number. Meanwhile $cz+d\to cx+d$, the smaller eigenvalue of $\gamma$ (at the repelling point
$\gamma'(x)=(cx+d)^{-2}>1$), which is positive because a positive word has trace greater than $2$. Hence $\epsilon=+1$ and $\psi(s+\ell)=\chi(\gamma)\psi(s)$, with
$\chi(\gamma)=\chi(L)^{\#L}\chi(R)^{\#R}=i^{\#L-\#R}$ because $\chi$ is a character.
\end{proof}

\begin{corollary}
The integral of $\psi$ over $n$ consecutive periods is $\sum_{k<n}i^{k(\#L-\#R)}$ times the one-period integral. So its
average over $n$ periods tends to $0$ as $n\to\infty$ unless $i^{\#L-\#R}=1$, that is $\#L\equiv\#R\pmod 4$. (The one-period
integral itself is in general not zero.)
\end{corollary}

The phase is a class function: $\chi$ is a character, so every positive word for the same conjugacy class gives the same
$i^{\#L-\#R}$. If the axis passes through a point of the orbit of $\omega$, where $D$ vanishes, $\psi$ has zeros there; the
theorem is unaffected.

The script \texttt{period\_phase.py} checks the theorem on eight words, at three starting points each, to $10^{-28}$.
On all $30\,674$ closed geodesics with trace at most $600$, a least-squares phase over two periods returns $i^{\#L-\#R}$ to
$1.1\cdot10^{-7}$~\cite{OriginAxiomRepo} (B1666).

\section{The surface average and the trace classes}\label{sec:avg}

The average of $f$ over the modular surface, $\frac3\pi\int_{\mathcal F}f\,\frac{dx\,dy}{y^2}$ over the standard fundamental domain,
is
\[
  \langle f\rangle=0.3150049063,
\]
the same at three Gauss--Legendre resolutions (\texttt{surface\_average.py}). By the equidistribution of closed
geodesics~\cite{Bowen1972,Margulis2004} and of their discriminant layers~\cite{Duke1988,Sarnak1982}, the average of $f$ along
the closed geodesics of a fixed trace $t$ tends to $\langle f\rangle$ as $t$ grows; all closed geodesics of trace $t$ have
the same length $2\operatorname{arccosh}(t/2)$, so the plain mean of their time averages is also the length-weighted one.
For the classes of trace at most $600$, the mean over the top quarter of traces ($451$ to $600$) is $0.3144$, against
$0.3692$ for traces at most $20$, with the distance falling roughly as $t^{-1/2}$~\cite{OriginAxiomRepo} (B1666). This is a
check of a computation against the theorem, not a new result.

\section*{Data, code and the use of AI}

The computations of Sections~\ref{sec:nil}--\ref{sec:avg} have scripts in \texttt{papers/P10\_theta\_triplet\_note/checks/}
of the public repository \url{https://github.com/originaxiom/origin-axiom}, and \texttt{verify.sh} runs them all. Two sets
of numbers are quoted from the record rather than recomputed here: the phase on the $30\,674$ closed geodesics of trace at
most $600$, and the trace-class means of Section~\ref{sec:avg} (both B1666). The computations and the drafting
were done with AI assistance (large-language-model sessions working in the repository under the author's direction);
the identities were checked symbolically or numerically as stated, and the author is responsible for the text.

\bibliographystyle{plain}
\bibliography{refs}
\end{document}
